\documentclass[11pt,oneside]{article}

\usepackage[utf-8]{inputenc}
\usepackage[margin=1in]{geometry}
\usepackage{amsmath}
\usepackage{amssymb}
\usepackage{amsthm}
\usepackage{mathtools}
\usepackage{booktabs}
\usepackage{array}
\usepackage{multirow}
\usepackage{xcolor}
\usepackage{listings}
\usepackage{xfrac}
\usepackage{hyperref}
\usepackage{fancyhdr}
\usepackage{tcolorbox}
\usepackage{mdframed}

% Code formatting
\lstset{
    language=Python,
    basicstyle=\ttfamily\small,
    keywordstyle=\color{blue},
    commentstyle=\color{gray},
    stringstyle=\color{red},
    breaklines=true,
    postbreak=\mbox{\textcolor{red}{$\hookrightarrow$}\space},
    showstringspaces=false,
    frame=single,
    backgroundcolor=\color{gray!10}
}

% Custom commands
\newcommand{\defn}[1]{\textbf{#1}}
\newcommand{\norm}[1]{\left\|#1\right\|}
\newcommand{\frobnorm}[1]{\left\|#1\right\|_F}
\newcommand{\inner}[2]{\langle #1, #2 \rangle}
\newcommand{\R}{\mathbb{R}}
\newcommand{\RR}[2]{\mathbb{R}^{#1 \times #2}}
\newcommand{\Real}{\mathbb{R}}
\newcommand{\N}{\mathbb{N}}
\newcommand{\grad}{\nabla}
\newcommand{\question}[1]{\textbf{Q: #1}}
\newcommand{\answer}{\textbf{A: }}

% Title and author
\title{\Large\textbf{Frequently Asked Questions (FAQ)}\\\Large Low-Rank Matrix Factorization Project}
\author{Comprehensive Documentation}
\date{\today}

% Header and footer
\pagestyle{fancy}
\fancyhf{}
\fancyhead[L]{Low-Rank Matrix Factorization FAQ}
\fancyhead[R]{\thepage}
\fancyfoot[C]{\small This document provides comprehensive documentation of the low-rank matrix factorization project.}

\hypersetup{
    colorlinks=true,
    linkcolor=blue,
    urlcolor=blue,
    citecolor=black,
    bookmarksnumbered=true,
}

\begin{document}

\maketitle

\begin{abstract}
This comprehensive FAQ document addresses the fundamental concepts, mathematical formulations, optimization algorithms, computational complexity, and practical considerations for low-rank matrix factorization. It serves as a complete reference guide for understanding the implementation, choosing appropriate methods, and troubleshooting common issues.
\end{abstract}

\tableofcontents
\newpage

% ============================================================================
\section{What is Low-Rank Matrix Factorization?}\label{sec:intro}
% ============================================================================

\question{What is Low-Rank Matrix Factorization?}

\answer Low-rank matrix factorization is a technique that decomposes a large matrix $M \in \Real^{m \times n}$ into a product of two smaller matrices:

\begin{equation}
    M \approx U V^T
\end{equation}

where $U \in \Real^{m \times r}$ and $V \in \Real^{n \times r}$, with $r \ll \min(m, n)$ being the \defn{rank} of the factorization.

\paragraph{Why is this useful?}

\begin{itemize}
    \item \textbf{Dimensionality reduction}: Compresses large matrices while preserving important structure
    \item \textbf{Noise removal}: Can recover underlying low-rank structure from noisy observations
    \item \textbf{Prediction}: Used in recommendation systems (e.g., MovieLens) to predict missing values
    \item \textbf{Computational efficiency}: Requires only $r(m + n)$ parameters instead of $mn$ for the full matrix
\end{itemize}

% ============================================================================
\section{Problem Formulation and Mathematical Framework}\label{sec:formulation}
% ============================================================================

\question{What is the exact optimization problem being solved?}

\answer We solve the \defn{regularized low-rank matrix factorization problem}:

\begin{equation}\label{eq:objective}
    \min_{U, V} f(U, V) = \frac{1}{2}\left\|UV^T - M\right\|_F^2 + \frac{\lambda}{2}\left(\|U\|_F^2 + \|V\|_F^2\right)
\end{equation}

where:
\begin{itemize}
    \item $U \in \Real^{m \times r}$ and $V \in \Real^{n \times r}$ are the \defn{factor matrices}
    \item $M \in \Real^{m \times n}$ is the observed matrix (dense or sparse)
    \item $\|\cdot\|_F$ denotes the \defn{Frobenius norm}: $\|A\|_F = \sqrt{\sum_{i,j} a_{ij}^2}$
    \item $\lambda \geq 0$ is the \defn{regularization parameter} (controls model complexity)
\end{itemize}

\medskip

\question{What about the sparse version?}

\answer For sparse matrices with observed entries $\{(i, j, M_{ij}) : (i,j) \in \Omega\}$, we solve:

\begin{equation}\label{eq:sparse_objective}
    \min_{U, V} f(U, V) = \frac{1}{2}\sum_{(i,j) \in \Omega}(U_i V_j^T - M_{ij})^2 + \frac{\lambda}{2}\left(\|U\|_F^2 + \|V\|_F^2\right)
\end{equation}

This optimizes only over observed ratings (e.g., in recommendation systems) rather than the entire matrix.

% ============================================================================
\section{Key Mathematical Symbols and Notation}\label{sec:notation}
% ============================================================================

\question{What are the standard mathematical notations used in this project?}

\answer Here are the key symbols and their meanings:

\begin{table}[h]
\centering
\begin{tabular}{lll}
\toprule
\textbf{Symbol} & \textbf{Meaning} & \textbf{Type/Domain} \\
\midrule
$U$ & First factor matrix & $\Real^{m \times r}$ \\
$V$ & Second factor matrix & $\Real^{n \times r}$ \\
$M$ & Target/observed matrix & $\Real^{m \times n}$ \\
$r$ & Rank (latent dimensionality) & scalar $\ll \min(m,n)$ \\
$m, n$ & Dimensions of $M$ & scalars \\
$\lambda$ & Regularization parameter & $\lambda > 0$ \\
$\grad_U f$ & Gradient w.r.t.\ $U$ & $\Real^{m \times r}$ \\
$\grad_V f$ & Gradient w.r.t.\ $V$ & $\Real^{n \times r}$ \\
$\alpha$ or $\eta$ & Step size/learning rate & scalar \\
$f(U, V)$ & Objective function value & $\Real$ \\
$\Omega$ & Set of observed entries (sparse) & index set \\
$\|\cdot\|_F$ & Frobenius norm & $\Real \to \Real_{\geq 0}$ \\
\bottomrule
\end{tabular}
\end{table}

% ============================================================================
\section{Optimization Methods: Algorithms and Approaches}\label{sec:methods}
% ============================================================================

\question{What three optimization methods does the project implement?}

\answer This project compares three state-of-the-art optimization algorithms.

% ----
\subsection{Gradient Descent with Armijo Backtracking (GD)}\label{subsec:gd}
% ----

\paragraph{Algorithm Overview}

At each iteration $k$:
\begin{enumerate}
    \item Compute gradient: $\mathbf{g}_k = \grad f(U_k, V_k) = [\grad_U f, \grad_V f]^T$
    \item Perform line search to find step size $\alpha_k$
    \item Update: $\theta_{k+1} = \theta_k - \alpha_k \mathbf{g}_k$
\end{enumerate}

\paragraph{Gradient Formulation}

For dense problems:
\begin{align}
    \grad_U f &= (UV^T - M)V + \lambda U \\
    \grad_V f &= (UV^T - M)^T U + \lambda V
\end{align}

For sparse problems:
\begin{align}
    \grad_U f &= \sum_{(i,j) \in \Omega} (U_i V_j^T - M_{ij})V_j + \lambda U \\
    \grad_V f &= \sum_{(i,j) \in \Omega} (U_i V_j^T - M_{ij})U_i + \lambda V
\end{align}

\paragraph{Armijo Line Search Condition}

We accept step size $\alpha_k$ if:
\begin{equation}
    f(\theta_{k+1}) \leq f(\theta_k) - c \alpha_k \|\mathbf{g}_k\|^2
\end{equation}

where $c = 10^{-4}$ (default) is the \defn{Armijo constant}. This ensures sufficient decrease in the objective function.

\paragraph{Backtracking Strategy}

\begin{itemize}
    \item Start with initial guess $\alpha = 1.0$
    \item If condition fails, shrink: $\alpha \leftarrow 0.5 \cdot \alpha$ (backtracking with $\rho = 0.5$)
    \item Continue until line search succeeds or $\alpha < 10^{-12}$ (failure)
\end{itemize}

\paragraph{Convergence Criterion}

Stops when gradient norm $\|\mathbf{g}_k\|_2 \leq 10^{-6}$ or max iterations reached.

% ----
\subsection{L-BFGS (Limited-memory BFGS)}\label{subsec:lbfgs}
% ----

\paragraph{Algorithm Overview}

L-BFGS is a \defn{quasi-Newton method} that approximates the second-order Hessian information without storing the full matrix.

\paragraph{Key Idea}

Instead of storing the full approximate Hessian $\mathbf{B}_k$ (requires $O(d^2)$ memory where $d = mr + nr$), L-BFGS stores only the last $m$ pairs of gradient differences and update vectors.

\paragraph{Update Rule}

\begin{equation}
    \theta_{k+1} = \theta_k - \alpha_k \mathbf{B}_k^{-1} \grad f(\theta_k)
\end{equation}

where $\mathbf{B}_k^{-1}$ is implicitly reconstructed from stored history.

\paragraph{Why Use L-BFGS Over GD?}

\begin{itemize}
    \item \textbf{Faster convergence}: Superlinear convergence vs.\ linear for GD
    \item \textbf{Adaptive step sizes}: Incorporates second-order information
    \item \textbf{Memory efficient}: $O(mr + nr)$ memory instead of $O((mr + nr)^2)$
\end{itemize}

\paragraph{Implementation}

We use \texttt{scipy.optimize.minimize} with \texttt{method='L-BFGS-B'} for box-constrained optimization.

% ----
\subsection{Alternating Least Squares (ALS)}\label{subsec:als}
% ----

\paragraph{Algorithm Overview}

ALS exploits the \defn{bilinear structure} of the problem by optimizing $U$ and $V$ separately in alternating fashion.

\paragraph{ALS Updates}

\textit{Fix $V$, optimize $U$:}
\begin{equation}
    U^* = \arg\min_U \frac{1}{2}\|UV^T - M\|_F^2 + \frac{\lambda}{2}\|U\|_F^2
\end{equation}

This is a least-squares problem with closed-form solution:
\begin{equation}
    U_i^* = (V^TV + \lambda I)^{-1} V^T M_i^T
\end{equation}

Or in matrix form:
\begin{equation}
    U^* = M(V + \lambda I)^{-1}
\end{equation}

\textit{Fix $U$, optimize $V$:}
\begin{equation}
    V^* = \arg\min_V \frac{1}{2}\|UV^T - M\|_F^2 + \frac{\lambda}{2}\|V\|_F^2
\end{equation}

Solution:
\begin{equation}
    V_j^* = (U^TU + \lambda I)^{-1} U^T M_j
\end{equation}

Or in matrix form:
\begin{equation}
    V^* = M^T(U + \lambda I)^{-1}
\end{equation}

\paragraph{Sparse Version}

For sparse data, we update row-by-row and column-by-column using only observed entries. For row $i$:
\begin{equation}
    U_i^* = (V_{\text{obs}}^T V_{\text{obs}} + \lambda I)^{-1} V_{\text{obs}}^T M_{i,\text{obs}}
\end{equation}

where $V_{\text{obs}}$ contains only columns corresponding to observed entries in row $i$.

\paragraph{Advantages}

\begin{itemize}
    \item \textbf{Exact solutions}: Each subproblem is solved exactly (or via pseudo-inverse)
    \item \textbf{Stable convergence}: No step size tuning needed
    \item \textbf{Highly efficient for sparse data}: Complexity depends on sparsity pattern
\end{itemize}

% ============================================================================
\section{Why These Optimization Techniques?}\label{sec:why}
% ============================================================================

\question{Why use gradient descent instead of directly solving the normal equations?}

\answer While the normal equations $\mathbf{A}^T\mathbf{A}\mathbf{x} = \mathbf{A}^Tb$ have closed-form solutions, directly solving them is problematic:

\begin{enumerate}
    \item \textbf{Non-convex problem}: The factorization problem is \defn{non-convex} in $(U, V)$ jointly. There are no global optimality conditions from standard linear algebra.
    \item \textbf{Numerical stability}: For ill-conditioned systems, solving normal equations amplifies numerical errors (condition number is squared).
    \item \textbf{Sparse data handling}: Gradient-based methods naturally handle sparse, high-dimensional problems.
    \item \textbf{Local convergence guarantees}: Gradient descent with backtracking ensures convergence to stationary points.
\end{enumerate}

\medskip

\question{Why use Armijo line search in gradient descent?}

\answer The Armijo condition provides:
\begin{itemize}
    \item \textbf{Guaranteed descent}: Each step reduces the objective by at least $c\alpha_k \|\mathbf{g}_k\|^2$
    \item \textbf{Adaptive step sizes}: Large steps when far from optimum, small steps when close
    \item \textbf{Efficiency}: Avoids expensive exact line search while maintaining convergence
    \item \textbf{Robustness}: Works without problem-specific knowledge
\end{itemize}

\medskip

\question{Why use L-BFGS?}

\answer L-BFGS combines advantages of both GD and Newton methods:
\begin{itemize}
    \item \textbf{Superlinear convergence}: Faster than GD, approaching Newton's quadratic convergence near optimum
    \item \textbf{Memory efficient}: Practical for high-dimensional problems $(mr + nr)$ parameters
    \item \textbf{Automatic step sizing}: Uses approximate Hessian information
    \item \textbf{Black-box solver}: No algorithm tuning needed (works well out-of-the-box)
\end{itemize}

\medskip

\question{When should we use ALS instead of GD or L-BFGS?}

\answer Use ALS when:
\begin{itemize}
    \item \textbf{Bilinear structure is critical}: Problem must be solved repeatedly with changing data
    \item \textbf{Sparse matrices dominate}: Row/column-wise updates exploit sparsity patterns
    \item \textbf{Computational efficiency matters}: ALS avoids expensive gradient computations per iteration
    \item \textbf{Convergence behavior is stable}: No step size tuning; often faster initial convergence
\end{itemize}

% ============================================================================
\section{Computational Complexity and Convergence Analysis}\label{sec:complexity}
% ============================================================================

\question{What is the computational complexity per iteration for each method?}

\answer

\begin{table}[h]
\centering
\small
\begin{tabular}{lllp{3cm}}
\toprule
\textbf{Method} & \textbf{Dense} & \textbf{Sparse} & \textbf{Comments} \\
\midrule
GD & $O(mnr)$ & $O(\text{nnz} \cdot r)$ & Gradient: $(UV^T - M)V$ \\
L-BFGS & $O(mnr)$ & $O(\text{nnz} \cdot r)$ & Like GD; Hessian: $O(mr)$ \\
ALS (dense) & $O(r^3 + mnr)$ & $O(r^3 + \text{nnz} \cdot r)$ & Solve: $O(r^3)$ \\
ALS (sparse) & --- & $O(\text{nnz} \cdot r^2/r)$ & Row/col updates only \\
\bottomrule
\end{tabular}
\caption{Per-iteration computational complexity. Here $\text{nnz}$ is the number of nonzero entries.}
\end{table}

\medskip

\question{What are the convergence rates?}

\answer

\paragraph{Gradient Descent}

\begin{itemize}
    \item \textbf{Local convergence rate}: Linear (geometric)
    \item \textbf{Rate constant}: Depends on Lipschitz constant $L$ of gradient
    \item \textbf{Asymptotic bound}: $\|\mathbf{g}_k\| = O(\rho^k)$ where $\rho < 1$
\end{itemize}

\paragraph{L-BFGS}

\begin{itemize}
    \item \textbf{Superlinear convergence}: $\|\mathbf{g}_k\| = o(\rho^k)$ for any $\rho > 0$
    \item \textbf{Typically faster}: Quadratic-like near optimum
    \item \textbf{No exact rate}: Depends on problem conditioning
\end{itemize}

\paragraph{ALS}

\begin{itemize}
    \item \textbf{Linear convergence}: Similar to GD but often with better constants
    \item \textbf{Problem-dependent}: Depends on spectrum of $(V^TV + \lambda I)^{-1}(V^TV)$ type matrices
    \item \textbf{Fast initial convergence}: Often beats GD for first 5--10 iterations
\end{itemize}

\medskip

\question{How do we measure convergence in practice?}

\answer We monitor:

\begin{enumerate}
    \item \textbf{Objective function}: $f(U_k, V_k)$ should monotonically decrease
    \item \textbf{Gradient norm}: 
    \begin{equation}
        \|\grad f\|_F = \sqrt{\|\grad_U f\|_F^2 + \|\grad_V f\|_F^2}
    \end{equation}
    Converged when $\|\grad f\|_F \leq 10^{-6}$
    
    \item \textbf{Relative change}:
    \begin{equation}
        \frac{|f(U_k, V_k) - f(U_{k-1}, V_{k-1})|}{|f(U_{k-1}, V_{k-1})|}
    \end{equation}
    Converged when this is $\leq 10^{-8}$
    
    \item \textbf{Wall-clock time}: Total execution time in seconds
\end{enumerate}

% ============================================================================
\section{Regularization and Its Role}\label{sec:regularization}
% ============================================================================

\question{Why do we use $L_2$ regularization with parameter $\lambda$?}

\answer Regularization serves multiple purposes:

\begin{enumerate}
    \item \textbf{Prevents overfitting}: 
    \begin{itemize}
        \item Without regularization, factors $U, V$ can grow arbitrarily large while maintaining small loss
        \item Regularization term $\lambda(\|U\|_F^2 + \|V\|_F^2)$ keeps them bounded
    \end{itemize}
    
    \item \textbf{Improves numerical stability}:
    \begin{itemize}
        \item Makes the problem better conditioned
        \item Helps with sparse data where some rows/columns have few observations
    \end{itemize}
    
    \item \textbf{Implicit smoothing}:
    \begin{itemize}
        \item Encourages simpler, smoother solutions
        \item Relevant for noisy data
    \end{itemize}
    
    \item \textbf{Connection to Bayesian interpretation}:
    \begin{itemize}
        \item Equivalent to $L_2$ prior on $U, V$: $P(U,V) \propto \exp(-\lambda(\|U\|^2 + \|V\|^2))$
    \end{itemize}
\end{enumerate}

\medskip

\question{How do we choose $\lambda$?}

\answer Common strategies:

\begin{enumerate}
    \item \textbf{Cross-validation}: Split data into train/validation, choose $\lambda$ minimizing validation error
    \item \textbf{Grid search}: Try $\lambda \in \{10^{-3}, 10^{-2}, 10^{-1}, 1.0\}$ and compare
    \item \textbf{Problem-specific}: 
    \begin{itemize}
        \item MovieLens: typically $\lambda = 10^{-2}$ works well
        \item Superconductivity: may need larger $\lambda$ due to noise
    \end{itemize}
    \item \textbf{Generalized cross-validation}: Automatic selection without explicit validation set
\end{enumerate}

% ============================================================================
\section{Initialization Strategies}\label{sec:initialization}
% ============================================================================

\question{How do we initialize $U_0$ and $V_0$?}

\answer We provide two strategies.

\subsection*{Random Initialization}

\begin{equation}
    U_0 \sim \mathcal{N}(0, \sigma^2 I), \quad V_0 \sim \mathcal{N}(0, \sigma^2 I)
\end{equation}

where $\sigma = 0.1$ (small scale to avoid large initial gradients).

\paragraph{Pros:} Simple, requires no prior information

\paragraph{Cons:} May converge to poor local minima

\subsection*{SVD-based Initialization}

\begin{enumerate}
    \item Compute truncated SVD: $M = U_{\text{SVD}}\Sigma_r V_{\text{SVD}}^T$ (rank-$r$ approximation)
    \item Initialize:
    \begin{equation}
        U_0 = U_{\text{SVD}} \sqrt{\Sigma_r}, \quad V_0 = V_{\text{SVD}} \sqrt{\Sigma_r}
    \end{equation}
    where $\sqrt{\Sigma_r}$ is element-wise square root.
\end{enumerate}

\paragraph{Pros}
\begin{itemize}
    \item Starts very close to optimal solution
    \item Better local optimum quality
    \item Faster convergence
\end{itemize}

\paragraph{Cons}
\begin{itemize}
    \item Expensive for large matrices ($O(mn \min(m,n))$)
    \item Requires dense matrix in memory
    \item May overfit for sparse data
\end{itemize}

% ============================================================================
\section{Handling Dense vs.\ Sparse Data}\label{sec:sparse}
% ============================================================================

\question{How does the implementation handle sparse matrices?}

\answer The project supports both \defn{dense} and \defn{sparse} representations.

\subsection*{Dense Problems}

\begin{itemize}
    \item Store full $m \times n$ matrix in memory
    \item Compute gradients: $\grad_U f = (UV^T - M)V + \lambda U$
    \item All matrix operations use standard linear algebra
\end{itemize}

\subsection*{Sparse Problems}

\begin{itemize}
    \item Store only nonzero entries as coordinate format: $(i, j, M_{ij})$
    \item Build \defn{lookup tables}:
    \begin{itemize}
        \item \texttt{row\_to\_cols[i]}: columns with nonzero entries in row $i$
        \item \texttt{row\_to\_vals[i]}: corresponding values
        \item \texttt{col\_to\_rows[j]}: rows with nonzero entries in column $j$
        \item \texttt{col\_to\_vals[j]}: corresponding values
    \end{itemize}
\end{itemize}

\subsection*{Gradient Computation (Sparse)}

\begin{equation}
    \grad_{U_i} f = \sum_{j \in \text{row\_to\_cols}[i]} (U_i V_j^T - M_{ij})V_j + \lambda U_i
\end{equation}

Only computes contributions from observed entries.

\subsection*{ALS for Sparse}

For each row $i$:
\begin{equation}
    U_i = (V_{\text{obs}}^T V_{\text{obs}} + \lambda I)^{-1} V_{\text{obs}}^T M_{i,\text{obs}}
\end{equation}

where $V_{\text{obs}} = V[\text{row\_to\_cols}[i]]$ contains only relevant columns.

\paragraph{Benefits}

\begin{itemize}
    \item Memory: $O(\text{nnz})$ instead of $O(mn)$
    \item Time: $O(\text{nnz} \cdot r)$ instead of $O(mnr)$
    \item Critical for massive datasets (MovieLens: $\sim$100K ratings, but $943 \times 1682 \approx 1.6\text{M}$ entries)
\end{itemize}

% ============================================================================
\section{Datasets and Evaluation Metrics}\label{sec:datasets}
% ============================================================================

\question{What datasets are used, and how do we evaluate performance?}

\answer The project tests on three dataset types.

\subsection*{1. Synthetic Low-Rank Matrices}

\paragraph{Generation}

\begin{itemize}
    \item Create true factors $U_* \in \Real^{m \times r}, V_* \in \Real^{n \times r}$
    \item Generate data: $M = U_* V_*^T + \text{noise}$, where noise $\sim \mathcal{N}(0, \sigma^2)$
    \item Default: $m=500, n=200, r=5, \sigma=0.1$
\end{itemize}

\paragraph{Evaluation}

\begin{equation}
    \text{Reconstruction Error} = \|U_{\text{recovered}}V_{\text{recovered}}^T - M_{\text{original}}\|_F
\end{equation}

\subsection*{2. Superconductivity Dataset}

\paragraph{Source:} UCI Machine Learning Repository (Dataset ID: 464)

\paragraph{Properties}
\begin{itemize}
    \item 21,263 samples $\times$ 81 features
    \item Regression problem (predicting superconducting temperature)
    \item Treated as dense matrix for factorization
\end{itemize}

\paragraph{Evaluation:} Test on held-out features/samples using reconstruction error

\subsection*{3. MovieLens 100K}

\paragraph{Source:} Grouplens MovieLens dataset

\paragraph{Properties}
\begin{itemize}
    \item 943 users $\times$ 1,682 movies
    \item $\sim$100,000 ratings (5-star scale)
    \item Extremely sparse: $\approx 6.3\%$ density
    \item Multiple train/test splits (u1.base/u1.test, etc.)
\end{itemize}

\paragraph{Evaluation Metric --- RMSE on Test Set}

\begin{equation}
    \text{RMSE}_{\text{test}} = \sqrt{\frac{1}{|\Omega_{\text{test}}|}\sum_{(i,j) \in \Omega_{\text{test}}}(U_i V_j^T - M_{ij})^2}
\end{equation}

where $\Omega_{\text{test}}$ are the held-out test ratings.

% ============================================================================
\section{Experiments and Comparisons}\label{sec:experiments}
% ============================================================================

\question{What experiments does the project run?}

\answer

\subsection*{Experiment 1: Method Comparison}

\textbf{Purpose:} Compare GD, L-BFGS, and ALS on standard setting

\textbf{Variables:} Fixed rank, fixed $\lambda$, vary dataset

\textbf{Metric:} Time to convergence, final objective value, solution quality

\subsection*{Experiment 2: Initialization Sensitivity}

\textbf{Purpose:} Assess random vs.\ SVD initialization impact

\textbf{Variables:} Initialization method, vary dataset

\textbf{Finding:} SVD init typically gives better solutions with faster convergence

\subsection*{Experiment 3: Rank Sensitivity}

\textbf{Purpose:} Study effect of chosen rank $r$ vs.\ true rank $r_{\text{true}}$

\textbf{Variables:} Range of $r$ values, fixed $\lambda$

\textbf{Findings}
\begin{itemize}
    \item Best performance at $r = r_{\text{true}}$
    \item Underfitting when $r < r_{\text{true}}$ (bias)
    \item Overfitting when $r > r_{\text{true}}$ (variance)
\end{itemize}

\subsection*{Experiment 4: Regularization Sensitivity}

\textbf{Purpose:} Optimize regularization parameter $\lambda$

\textbf{Variables:} Grid of $\lambda$ values

\textbf{Findings:} U-shaped error curve; optimal $\lambda$ balances bias-variance

\subsection*{Experiment 5: Loss Landscape}

\textbf{Purpose:} Visualize optimization landscape

\textbf{Exploration:} 2D slice through factor space

\textbf{Findings:} Multiple local minima; algorithms converge to different solutions from different initializations

% ============================================================================
\section{Stopping Criteria and Convergence Diagnostics}\label{sec:stopping}
% ============================================================================

\question{When does the algorithm stop iterating?}

\answer Algorithms stop when ANY of these conditions is met:

\begin{enumerate}
    \item \textbf{Gradient convergence}:
    \begin{equation}
        \|\grad f(U_k, V_k)\|_F \leq 10^{-6}
    \end{equation}
    
    \item \textbf{Objective stalling}:
    \begin{equation}
        \frac{|f_k - f_{k-1}|}{|f_{k-1}|} \leq 10^{-8}
    \end{equation}
    (Relative change too small to matter)
    
    \item \textbf{Maximum iterations reached}: Default: 500 iterations
    
    \item \textbf{Failure modes}:
    \begin{itemize}
        \item Line search fails (GD only): Cannot find acceptable step size
        \item Singular linear system (ALS only): Pseudo-inverse fallback used
    \end{itemize}
\end{enumerate}

\paragraph{Status Field}

The history dictionary tracks termination reason:
\begin{itemize}
    \item \texttt{"converged"}: Gradient converged to tolerance
    \item \texttt{"objective\_stalled"}: Function value stopped improving
    \item \texttt{"max\_iter"}: Hit iteration limit
    \item \texttt{"line\_search\_failed"}: GD backtracking failed
\end{itemize}

% ============================================================================
\section{Implementation Details and Numerical Considerations}\label{sec:implementation}
% ============================================================================

\question{What numerical precautions are taken?}

\answer

\begin{enumerate}
    \item \textbf{Pseudo-inverse fallback (ALS)}:
    \begin{itemize}
        \item Try to solve $\mathbf{A}x = b$ via \texttt{np.linalg.solve()}
        \item If matrix is singular, use \texttt{np.linalg.pinv()} (pseudo-inverse)
        \item Ensures robust behavior even with ill-conditioned systems
    \end{itemize}
    
    \item \textbf{Gradient norm computation}:
    \begin{equation}
        \|\grad f\|_2 = \sqrt{\sum_{i,j}(\grad_U f)_{ij}^2 + (\grad_V f)_{ij}^2}
    \end{equation}
    Computed efficiently using summed squared elements
    
    \item \textbf{Step size bounds}:
    \begin{itemize}
        \item Minimum: $\alpha_{\min} = 10^{-12}$ (GD backtracking)
        \item Initial guess: $\alpha_0 = 1.0$
        \item Never exceeds: $\alpha_{\max} = 2 \times \alpha_{\text{previous}}$
    \end{itemize}
    
    \item \textbf{Packing/unpacking factors}:
    \begin{itemize}
        \item For optimization algorithms requiring vector input (L-BFGS), concatenate:
        \begin{equation}
            \mathbf{x} = [U_{\text{vec}}; V_{\text{vec}}]
        \end{equation}
        \item Allows standard optimizers to work with matrix problems
    \end{itemize}
\end{enumerate}

% ============================================================================
\section{Performance Benchmarking and Timing}\label{sec:benchmarking}
% ============================================================================

\question{How is computational time measured?}

\answer Using \texttt{time.perf\_counter()} for wall-clock time:

\begin{lstlisting}
start = time.perf_counter()
# ... run optimization ...
elapsed = time.perf_counter() - start
\end{lstlisting}

\medskip

\textbf{Typical timings} (moderate hardware, synthetic data $500 \times 200$, rank 5):

\begin{table}[h]
\centering
\begin{tabular}{llll}
\toprule
\textbf{Method} & \textbf{Iterations} & \textbf{Time} & \textbf{Time/Iteration} \\
\midrule
GD & 150--300 & 2--5s & $\sim$15ms \\
L-BFGS & 40--80 & 1--3s & $\sim$30ms \\
ALS & 20--50 & 0.5--2s & $\sim$25ms \\
\bottomrule
\end{tabular}
\end{table}

\medskip

\question{What affects convergence speed?}

\answer

\begin{enumerate}
    \item \textbf{Initialization quality}
    \begin{itemize}
        \item SVD init: 5--10$\times$ faster than random
    \end{itemize}
    
    \item \textbf{Rank selection}
    \begin{itemize}
        \item Too large: Overfitting, slower convergence
        \item Too small: Underfitting, may not converge to good solution
    \end{itemize}
    
    \item \textbf{Regularization strength $\lambda$}
    \begin{itemize}
        \item Small $\lambda$: Longer to converge, higher risk of overfitting
        \item Large $\lambda$: Faster, but solution less accurate
    \end{itemize}
    
    \item \textbf{Data properties}
    \begin{itemize}
        \item Sparsity: Sparse data usually converges slower
        \item Noise: Higher noise requires more iterations
    \end{itemize}
    
    \item \textbf{Problem conditioning}
    \begin{itemize}
        \item Ill-conditioned: All methods slow down
        \item Well-conditioned: Superlinear methods (L-BFGS) excel
    \end{itemize}
\end{enumerate}

% ============================================================================
\section{Practical Tips and Troubleshooting}\label{sec:troubleshooting}
% ============================================================================

\question{Which method should I use for my problem?}

\answer

\begin{table}[h]
\centering
\begin{tabular}{lll}
\toprule
\textbf{Problem Type} & \textbf{Recommended} & \textbf{Reason} \\
\midrule
Dense, small--medium & L-BFGS & Best all-around \\
Dense, very large & GD & Memory efficient \\
Sparse, large & ALS & Exploits sparsity \\
Quick prototyping & GD & Simple, reliable \\
Production system & ALS or L-BFGS & Stable, predictable \\
\bottomrule
\end{tabular}
\end{table}

\medskip

\question{How do I avoid poor local minima?}

\answer

\begin{enumerate}
    \item \textbf{Use SVD initialization} (if computationally feasible)
    \item \textbf{Try multiple random seeds} and pick best result
    \item \textbf{Increase rank slightly} (may help escape bad minima)
    \item \textbf{Reduce regularization} (but watch for overfitting)
    \item \textbf{Increase iterations} (sometimes helps find better solution)
\end{enumerate}

\medskip

\question{My solution has high reconstruction error. What should I do?}

\answer

\textbf{Check these issues}:

\begin{enumerate}
    \item \textbf{Is rank sufficient?} Try $r + 1$ or $r + 2$
    \item \textbf{Is $\lambda$ too large?} Reduces model capacity
    \item \textbf{Is data noisy?} May need larger $\lambda$ or accept higher error
    \item \textbf{Did algorithm converge?} Check gradient norm and objective value
    \item \textbf{Is initialization poor?} Try SVD or multiple random seeds
\end{enumerate}

\medskip

\question{Convergence is very slow. How can I speed it up?}

\answer

\begin{enumerate}
    \item \textbf{Switch to L-BFGS} (superlinear vs.\ linear GD)
    \item \textbf{Use SVD initialization} (much better starting point)
    \item \textbf{Reduce tolerance} (converge to $10^{-5}$ instead of $10^{-6}$)
    \item \textbf{Optimize $\lambda$} (may be poorly chosen)
    \item \textbf{Check data quality} (noisy data requires more iterations)
    \item \textbf{For sparse data, use ALS} (best for MovieLens-type problems)
\end{enumerate}

% ============================================================================
\section{Mathematical Background References}\label{sec:references}
% ============================================================================

\question{Where can I learn more about the underlying mathematics?}

\answer

\textbf{Key Concepts}:

\begin{table}[h]
\centering
\small
\begin{tabular}{lll}
\toprule
\textbf{Topic} & \textbf{Relevance} & \textbf{Key Papers/Books} \\
\midrule
Matrix factorization & Core & Eckart \& Young (1936); Lee \& Seung (2001) \\
Optimization & Methods & Boyd \& Vandenberghe (2004) Ch.\ 1--3 \\
Gradient descent & GD method & Nesterov (2003) \\
Quasi-Newton methods & L-BFGS & Dennis \& Moré (1977); Nocedal \& Wright (2006) \\
Alternating optimization & ALS & Csiszár \& Tusnády (1984); Razaviyayn \textit{et al.} (2013) \\
Line search & Armijo & Armijo (1966); Nocedal \& Wright Ch.\ 3 \\
\bottomrule
\end{tabular}
\end{table}

% ============================================================================
\section{Common Questions About Results and Outputs}\label{sec:outputs}
% ============================================================================

\question{What does the history dictionary contain?}

\answer

\begin{lstlisting}
history = {
    "method": str,                  # "GD", "L-BFGS", or "ALS"
    "obj": List[float],             # Objective values over iterations
    "grad_norm": List[float],       # Gradient norms over iterations
    "time": List[float],            # Cumulative time in seconds
    "step_size": List[float],       # Step sizes (GD only)
    "U_final": ndarray,             # Final U matrix
    "V_final": ndarray,             # Final V matrix
    "status": str,                  # Convergence status
    "iterations": int,              # Total iterations taken
}
\end{lstlisting}

\medskip

\question{How do I extract and visualize results?}

\answer

\begin{lstlisting}
# Convergence plot
import matplotlib.pyplot as plt
for method, hist in results.items():
    plt.semilogy(hist["iterations"], 
                 hist["grad_norm"], 
                 label=method)
plt.legend()
plt.xlabel("Iteration")
plt.ylabel("Gradient Norm")
plt.show()

# Reconstruction error
U_final = history["U_final"]
V_final = history["V_final"]
reconstruction = U_final @ V_final.T
error = np.linalg.norm(reconstruction - 
                       original_matrix)
\end{lstlisting}

% ============================================================================
\section{Frequently Encountered Errors}\label{sec:errors}
% ============================================================================

\question{I get ``LinAlgError: Singular matrix'' during ALS. What should I do?}

\answer This happens when $V^T V + \lambda I$ or $U^T U + \lambda I$ is singular.

\textbf{Solutions}:

\begin{enumerate}
    \item \textbf{Increase $\lambda$}: Use larger regularization parameter (e.g., $10 \lambda$)
    \item \textbf{Algorithm fallback}: Code automatically uses pseudo-inverse (ensure \texttt{allow\_pinv=True})
    \item \textbf{Check rank}: May be too large for dataset size
    \item \textbf{Data validation}: Ensure no all-zero rows/columns in data
\end{enumerate}

\medskip

\question{Gradient check fails. What does this mean?}

\answer Gradient check compares analytical gradients with numerical finite-difference approximation:

\begin{equation}
    \text{relative\_error} = \frac{\|\grad f_{\text{analytical}} - \grad f_{\text{numerical}}\|}{\|\grad f_{\text{numerical}}\|}
\end{equation}

\textbf{If check fails} ($> 10^{-6}$ typically):

\begin{enumerate}
    \item \textbf{Bug in gradient code}: Review implementation
    \item \textbf{Numerical precision issue}: Use higher precision (float64 if using float32)
    \item \textbf{Step size too large}: Reduce for numerical differentiation
\end{enumerate}

% ============================================================================
\section{Relationship Between Methods}\label{sec:relationship}
% ============================================================================

\question{How do GD, L-BFGS, and ALS relate to each other?}

\answer

\vspace{0.5cm}

\begin{mdframed}[backgroundcolor=gray!10]
All three optimize:
\begin{equation*}
    \min_{U,V} \tfrac{1}{2}\|UV^T - M\|_F^2 + \lambda(\|U\|_F^2 + \|V\|_F^2)
\end{equation*}

\medskip

\textbf{GD}: Uses only first-order (gradient) information
\begin{itemize}
    \item Simple, stable
    \item Slow convergence (linear)
\end{itemize}

\textbf{L-BFGS}: Uses approximate second-order (Hessian) information
\begin{itemize}
    \item More complex, faster convergence (superlinear)
    \item Memory efficient (limited-memory variant)
\end{itemize}

\textbf{ALS}: Exploits bilinear structure; alternates optimization
\begin{itemize}
    \item Solves subproblems exactly via least-squares
    \item Natural for factorization; very efficient for sparse data
\end{itemize}
\end{mdframed}

\paragraph{Key Insight}

Different methods make different trade-offs between:
\begin{itemize}
    \item \textbf{Convergence speed} (iterations to convergence)
    \item \textbf{Per-iteration cost} (time per iteration)
    \item \textbf{Implementation complexity} (code simplicity)
    \item \textbf{Hyperparameter tuning} (algorithm parameters to set)
\end{itemize}

% ============================================================================
\section{FAQ Meta: How to Use This Document}\label{sec:meta}
% ============================================================================

\question{How should I use this FAQ?}

\answer

\begin{itemize}
    \item \textbf{For understanding concepts}: Start with Sections~\ref{sec:intro}--\ref{sec:notation} (Low-Rank, Problem, Notation)
    \item \textbf{For algorithm details}: Read Sections~\ref{sec:methods}--\ref{sec:complexity} (Methods, Why, Complexity)
    \item \textbf{For practical usage}: Jump to Sections~\ref{sec:implementation}--\ref{sec:errors} (Implementation, Troubleshooting)
    \item \textbf{For comparison}: Use tables in Sections~\ref{sec:methods}, \ref{sec:complexity}, and \ref{sec:troubleshooting}
    \item \textbf{For mathematical rigor}: Sections~\ref{sec:formulation}, \ref{sec:methods}, \ref{sec:complexity}, and \ref{sec:references} have detailed formulas
\end{itemize}

% ============================================================================
\section{Quick Reference Card}\label{sec:quick_ref}
% ============================================================================

\begin{table}[h]
\centering
\begin{tabular}{ll}
\toprule
\textbf{Question} & \textbf{Section} \\
\midrule
What is low-rank factorization? & \ref{sec:intro} \\
What problem are we solving? & \ref{sec:formulation}, \ref{sec:notation} \\
Which algorithm is fastest? & \ref{sec:methods}, \ref{sec:complexity}, \ref{sec:troubleshooting} \\
How do I choose hyperparameters? & \ref{sec:regularization}, \ref{sec:troubleshooting} \\
Why is my solution bad? & \ref{sec:troubleshooting} \\
What datasets are available? & \ref{sec:datasets} \\
How do I interpret results? & \ref{sec:outputs} \\
Something went wrong! & \ref{sec:errors} \\
\bottomrule
\end{tabular}
\caption{Quick reference guide with section cross-references.}
\end{table}

% ============================================================================
\appendix
% ============================================================================

\section{Mathematical Symbols Summary}\label{app:symbols}

\subsection*{Matrices and Vectors}

\begin{tabularx}{\textwidth}{llX}
\toprule
\textbf{Symbol} & \textbf{Dimension} & \textbf{Interpretation} \\
\midrule
$U$ & $m \times r$ & First low-rank factor \\
$V$ & $n \times r$ & Second low-rank factor \\
$M$ & $m \times n$ & Target matrix \\
$I$ & $r \times r$ & Identity matrix \\
$\mathbf{g}_k$ & $(mr+nr)$-d & Stacked gradient \\
\bottomrule
\end{tabularx}

\subsection*{Operators and Norms}

\begin{tabularx}{\textwidth}{llX}
\toprule
\textbf{Notation} & \textbf{Type} & \textbf{Definition} \\
\midrule
$\|\cdot\|_F$ & Frobenius norm & $\|A\|_F = \sqrt{\sum_{i,j} a_{ij}^2}$ \\
$\|\cdot\|_2$ & Euclidean norm & $\|v\|_2 = \sqrt{\sum_i v_i^2}$ \\
$\grad f$ & Gradient & Direction of steepest ascent \\
$\text{vec}(\cdot)$ & Vectorization & Stack matrix into vector \\
$^T$ & Transpose & Flip matrix rows/columns \\
$\otimes$ & Kronecker product & Tensor product of matrices \\
\bottomrule
\end{tabularx}

\subsection*{Frequently Used Acronyms}

\begin{tabularx}{\textwidth}{llX}
\toprule
\textbf{Acronym} & \textbf{Full Name} & \textbf{Usage} \\
\midrule
GD & Gradient Descent & Main first-order method \\
L-BFGS & Limited-memory BFGS & Quasi-Newton method \\
ALS & Alternating Least Squares & Bilinear optimization \\
SVD & Singular Value Decomposition & Initialization and analysis \\
RMSE & Root Mean Squared Error & Evaluation metric \\
nnz & Number of nonzeros & Sparse matrix property \\
\bottomrule
\end{tabularx}

% ============================================================================
\section{Key Equations Summary}\label{app:equations}
% ============================================================================

\paragraph{Objective Function (Dense)}
\begin{equation*}
    f(U, V) = \frac{1}{2}\|UV^T - M\|_F^2 + \frac{\lambda}{2}(\|U\|_F^2 + \|V\|_F^2)
\end{equation*}

\paragraph{Gradient Descent Update}
\begin{equation*}
    U_{k+1} = U_k - \alpha_k \grad_U f(U_k, V_k), \quad V_{k+1} = V_k - \alpha_k \grad_V f(U_k, V_k)
\end{equation*}

\paragraph{ALS U-Update}
\begin{equation*}
    U^* = M(V + \lambda I)^{-1}
\end{equation*}

\paragraph{ALS V-Update}
\begin{equation*}
    V^* = M^T(U + \lambda I)^{-1}
\end{equation*}

\paragraph{Armijo Condition}
\begin{equation*}
    f(\theta_{k+1}) \leq f(\theta_k) - c\alpha_k \|\grad f(\theta_k)\|^2
\end{equation*}

\newpage

\section{Further Reading}\label{app:reading}

\subsection*{Foundational Papers}

\begin{itemize}
    \item \textbf{Eckart, C., \& Young, G.} (1936). ``The approximation of one matrix by another of lower rank.'' \textit{Psychometrika}, 1(3), 211--218.
    
    \item \textbf{Nocedal, J., \& Wright, S. J.} (2006). \textit{Numerical Optimization} (2nd ed.). Springer.
    
    \item \textbf{Boyd, S., \& Vandenberghe, L.} (2004). \textit{Convex Optimization}. Cambridge University Press.
\end{itemize}

\subsection*{Method-Specific References}

\begin{itemize}
    \item \textbf{Nesterov, Y.} (2003). ``Introductory lectures on convex optimization.'' In \textit{Optimization Methods and Applications}. Kluwer Academic Publishers.
    
    \item \textbf{Dennis, J. E., \& Moré, J. J.} (1977). ``Quasi-Newton methods, motivation and theory.'' \textit{SIAM Review}, 19(1), 46--89.
    
    \item \textbf{Razaviyayn, M., Hong, M., \& Luo, Z. Q.} (2013). ``A unified convergence analysis of block successive minimization methods for nonsmooth optimization.'' \textit{SIAM Journal on Optimization}, 23(2), 733--756.
\end{itemize}

\subsection*{Application-Specific References}

\begin{itemize}
    \item \textbf{Lee, D. D., \& Seung, H. S.} (2001). ``Algorithms for non-negative matrix factorization.'' In \textit{Advances in Neural Information Processing Systems} (pp.\ 556--562).
    
    \item \textbf{Koren, Y., Bell, R., \& Volinsky, C.} (2009). ``Matrix factorization techniques for recommender systems.'' \textit{Computer}, 42(8), 30--37.
\end{itemize}

\end{document}
